\documentclass[12pt]{article}
\usepackage[utf8]{inputenc}
\usepackage[T1]{fontenc}
\usepackage{hyperref}
\title{Kuce Language Resources for K-12 Curriculum Development}
\author{Compiled by Claude Sonnet 5.5 (claude-sonnet-5-5) for Andr\'as Kornai}
\date{}
\begin{document}
\maketitle

\section{Introduction}
Kuce (also Ce, Che, Kuche, Rukuba, Inchazi; ISO 639-3 ruk) is a Plateau language of Plateau State, Nigeria. Wikipedia gives about 100,000 speakers (2003), and Joshua Project gives 183,000 people in the group. The people are mostly farmers. Search turned up very little, and nothing aimed at schools.

\section{Primary Educational Resources}
No primer, textbook or school programme for Kuce was found.

\section{Linguistic and Reference Materials}
Ce language (Wikipedia), a general overview: \url{https://en.wikipedia.org/wiki/Ce_language}

\section{Religious and Community Resources}
Joshua Project reports that no Bible translation exists. Church and community leaders want a New Testament, starting with Luke, Acts and James: \url{https://www.joshuaproject.net/index.php/languages/ruk}

\section{Digital and Audio Resources}
Global Recordings Network lists a Kuce page (audio recordings may be available there; I did not verify what it holds): \url{https://globalrecordings.net/en/language/8721}

\section{Curriculum Development Context}
Kuce has no known school use and little documentation online. A curriculum would have to start with community documentation: an orthography, a word list and recorded stories made with Kuce speakers. This survey found no existing teaching materials.
\end{document}